% Image credits for the photographs and AI illustrations of Book 1 (Primary and Middle School).
% Machine-readable license records live in images/book1/CREDITS.md; keep the
% two files in sync. The Book 1 writer owns this file: add one \item per
% photograph (inside the itemize, by chapter) as photographs land.
% Portuguese (Brazilian) edition: translated from frontmatter/image-credits.tex.
\chapter*{Créditos das imagens}

Os desenhos, os mapas e os gráficos deste livro são originais; os mapas e
os gráficos desenhados a partir de dados publicados citam a sua fonte na
legenda ou nesta lista. As fotografias são usadas sob as suas respectivas
licenças livres, com agradecimentos aos seus autores; os links completos
das fontes e os endereços das licenças de todas as fotografias estão
listados em \texttt{images/book1/CREDITS.md}, no repositório do livro.

\bigskip
Ao lado das fotografias e dos desenhos originais, algumas figuras são
ilustrações geradas por IA, produzidas com a geração de imagens da OpenAI
a partir de instruções escritas pelos editores do livro, e revisadas
quanto à exatidão científica antes de serem incluídas. Elas mostram
apenas cenas genéricas, nunca um lugar real, uma pessoa real ou uma
amostra real. A instrução completa de cada imagem gerada está registrada
em \texttt{images/book1/ai/PROMPTS.md}, no repositório.

\bigskip
\noindent Fotografias e dados, por capítulo:
\begin{itemize}\small
  \item Cap.~\ref{ch:g2:water-on-earth} (a água na Terra): a Terra vista da Apollo~17, 7~de~dezembro de 1972 (``The Blue Marble''); NASA, domínio público (via Wikimedia Commons). Os 100 quadrados são calculados a partir do relevo NOAA ETOPO1; os 100 baldes, a partir da tabela da U.S. Geological Survey Water Science School (segundo Shiklomanov, 1993).
  \item Cap.~\ref{ch:g3:stones-pebbles-sand} (pedras, seixos e areia): um muro de pedra seca de granito; fotografia de Tony Atkin, Geograph, CC BY-SA 2.0 (via Wikimedia Commons). Telhados de lajes de pedra na Suíça; fotografia de Michael C. Rygel, CC BY-SA 3.0 (via Wikimedia Commons).
  \item Cap.~\ref{ch:g3:continents-and-oceans} (o globo): os dois globos e o mapa-múndi são desenhados a partir dos litorais do Natural Earth (domínio público).
  \item Cap.~\ref{ch:g4:clouds-rain-snow} (nuvens, chuva e neve): cúmulos sobre o Kansas (S.~Corfidi), manchas de cirros (R.~F.~Kresge), uma camada de estratos em volta de picos de montanha, uma tempestade sobre o deserto do Arizona (coleção de B.~Hooke), uma pedra de granizo muito grande; todas da NOAA Photo Library, domínio público (via Wikimedia Commons).
  \item Cap.~\ref{ch:g4:volcanoes-earthquakes-first-look} (vulcões e terremotos): lava entrando no oceano em Kilauea, USGS Hawaiian Volcano Observatory; o monte Santa Helena, 18~de~maio de 1980, fotografia de Austin Post, USGS; uma casa arrancada da sua fundação pelo terremoto de Loma Prieta de 1989, fotografia de J.~K.~Nakata, USGS; todas em domínio público (via Wikimedia Commons). O mapa usa a lista de vulcões do Holoceno do Global Volcanism Program da Smithsonian e o catálogo de terremotos da U.S. Geological Survey.
  \item Cap.~\ref{ch:g4:world-physical-map} (o mapa físico): o Himalaia e a Amazônia fotografados da Estação Espacial Internacional, NASA, domínio público (via Wikimedia Commons). O mapa físico é desenhado a partir do relevo NOAA ETOPO1, com rios, lagos e litorais do Natural Earth (domínio público).
  \item Cap.~\ref{ch:g5:water-shapes-the-land} (a água modela a terra): o rio Mamoré fotografado da Estação Espacial Internacional (ISS007-E-10797) e o delta do Nilo visto pelo instrumento MODIS do satélite Terra, NASA, domínio público (via Wikimedia Commons).
  \item Cap.~\ref{ch:g5:weather-and-climate} (tempo e clima): o deserto de Sonora, U.S. National Park Service, Saguaro National Park, domínio público (via Wikimedia Commons). Os climogramas são desenhados a partir das NOAA NCEI U.S. Climate Normals 1991--2020 (domínio público).
  \item Cap.~\ref{ch:g5:natural-hazards} (as ameaças naturais): o olho do furacão Isabel fotografado da Estação Espacial Internacional por Ed Lu, NASA, domínio público; o deslizamento de terra de 2014 no estado de Washington, Mark Reid, U.S. Geological Survey, domínio público (ambos via Wikimedia Commons).
  \item Cap.~\ref{ch:g6:earth-a-planet} (a Terra, um planeta): Vênus vista pela Mariner~10 (NASA/JPL-Caltech); a Terra vista da Apollo~17 (NASA); \emph{Earthrise} visto da Apollo~8 (NASA, Bill Anders); Marte e Júpiter vistos pelo telescópio espacial Hubble (NASA, ESA, the Hubble Heritage Team e A.~Simon); todos em domínio público (via Wikimedia Commons). O \emph{pálido ponto azul} visto pela Voyager~1 em 1990, versão de 2020 (NASA/JPL-Caltech, PIA23645, domínio público). Dados dos planetas da ficha técnica planetária da NASA.
  \item Cap.~\ref{ch:g6:latitude-longitude} (latitude e longitude): o H4 de Harrison no National Maritime Museum, em Greenwich, fotografia de Colonel Warden, CC BY-SA 3.0 (via Wikimedia Commons). Litorais e cidades do Natural Earth (domínio público).
  \item Cap.~\ref{ch:g6:rocks-and-minerals} (rochas e minerais): granito, basalto, calcário fossilífero, arenito, mármore, folhelho, um cristal de quartzo e uma lâmina de mica muscovita, fotografias de James St.~John, CC BY 2.0; um telhado de ardósia, fotografia de Øyvind Holmstad, CC0 (todas via Wikimedia Commons).
  \item Cap.~\ref{ch:g6:inside-the-earth} (o interior da Terra): Inge Lehmann em 1932, fotografia de Even Neuhaus, Biblioteca Real da Dinamarca, domínio público (via Wikimedia Commons). As camadas, as densidades e as velocidades das ondas vêm do Preliminary Reference Earth Model (A.~M. Dziewonski e D.~L. Anderson, 1981).
  \item Cap.~\ref{ch:g6:water-underground} (a água subterrânea): Big Spring, no Missouri, fotografia de Johnston9494, CC BY-SA 4.0 (via Wikimedia Commons). Porosidades segundo R.~C. Heath, \emph{Basic Ground-Water Hydrology}, U.S. Geological Survey (1983); partes da água doce da USGS Water Science School.
  \item Cap.~\ref{ch:g7:plate-tectonics} (tectônica de placas): os continentes do sul reunidos, com as faixas de fósseis de \emph{Cynognathus}, \emph{Lystrosaurus}, \emph{Mesosaurus} e \emph{Glossopteris}; U.S. Geological Survey, \emph{This Dynamic Earth}, domínio público.
  \item Cap.~\ref{ch:g7:plate-tectonics}: a falha de San Andreas atravessando a planície de Carrizo; fotografia de Robert E. Wallace, U.S. Geological Survey, domínio público.
  \item Cap.~\ref{ch:g7:plate-tectonics}: limites de placas segundo P.~Bird (2003), GeoJSON de H.~Ahlenius (ODC-By 1.0); litorais do Natural Earth (domínio público); terremotos do catálogo da U.S. Geological Survey; vulcões do Global Volcanism Program da Smithsonian.
  \item Cap.~\ref{ch:g7:volcanoes} (os vulcões): o monte Santa Helena visto de Johnston Ridge em 17~de~maio e em 10~de~setembro de 1980, fotografias de Harry Glicken, U.S. Geological Survey; uma fonte de lava em Kilauea, 1983, e o Mauna Loa visto do ar, 1985, fotografias de J.~D. Griggs, U.S. Geological Survey; a coluna eruptiva do monte Pinatubo, 12~de~junho de 1991, fotografia de Dave Harlow, U.S. Geological Survey; todas em domínio público (via Wikimedia Commons). O mapa e as contagens usam a lista de vulcões do Holoceno do Global Volcanism Program da Smithsonian; limites de placas segundo P.~Bird (2003).
  \item Cap.~\ref{ch:g7:earthquakes} (os terremotos): uma aldeia perto do litoral de Sumatra depois do tsunami de 26~de~dezembro de 2004, fotografia da U.S. Navy por Philip A. McDaniel; as escarpas de falha do terremoto de Borah Peak de 1983, fotografia de R.~E. Wallace, U.S. Geological Survey; a estrada afundada no lago Hebgen, 1959, fotografia de I.~J. Witkind, U.S. Geological Survey; todas em domínio público (via Wikimedia Commons). As contagens e as profundidades usam o catálogo de terremotos da U.S. Geological Survey; a tabela de intensidades é abreviada da escala de Mercalli modificada da U.S. Geological Survey.
  \item Cap.~\ref{ch:g7:sediments-and-fossils} (sedimentos e fósseis): o Grand Canyon visto de Hermits Rest, fotografia do NPS por Kristen M. Caldon, domínio público; uma amonite, uma trilobita e marcas de ondulação no Capitol Reef National Park, fotografias de James St.~John, CC BY 2.0; pegadas de dinossauros em Dinosaur Ridge, fotografia de Dexter Perkins, CC0 (todas via Wikimedia Commons).
  \item Cap.~\ref{ch:g7:geological-time} (as idades da Terra): James Hutton, retrato de Henry Raeburn, domínio público (via Wikimedia Commons). As idades e as cores da escala do tempo são as da Carta Cronoestratigráfica Internacional (Comissão Internacional de Estratigrafia).
  \item Cap.~\ref{ch:g8:the-atmosphere} (a atmosfera): lançamento de um balão meteorológico em Topeka, no Kansas, fotografia de Stephen Corfidi, NOAA/National Weather Service; a borda da Terra ao pôr do sol vista da Estação Espacial Internacional, NASA, Expedição~23 da ISS; ambas em domínio público (via Wikimedia Commons). Os perfis são calculados a partir da \emph{U.S. Standard Atmosphere, 1976} (NOAA, NASA e U.S. Air Force).
  \item Cap.~\ref{ch:g8:winds-global-circulation} (os ventos e a circulação geral): as Américas vistas pelo satélite GOES-16 em 26~de~julho de 2022, NOAA, domínio público (via Wikimedia Commons). Os litorais do mapa dos alísios são do Natural Earth (domínio público).
  \item Cap.~\ref{ch:g8:fronts-and-weather-maps} (massas de ar, frentes e mapas do tempo): uma depressão ao largo da Islândia vista pelo satélite Terra, NASA/GSFC, MODIS Rapid Response Team, Jacques Descloitres, domínio público (via Wikimedia Commons). Litorais do mapa das massas de ar do Natural Earth (domínio público).
  \item Cap.~\ref{ch:g8:the-oceans} (os oceanos): o HMS \emph{Challenger} preparando uma sondagem, 1872, gravura dos relatórios da expedição do Challenger reproduzida em W.~A. Herdman, \emph{Founders of Oceanography} (1923); a Corrente do Golfo vista pelo instrumento MODIS do satélite Aqua, NASA Earth Observatory; ambas em domínio público (via Wikimedia Commons). Os perfis do fundo do oceano são desenhados a partir do relevo NOAA ETOPO1; os litorais dos mapas dos giros e das boias são do Natural Earth (domínio público).
  \item Cap.~\ref{ch:g8:ice-on-earth} (o gelo na Terra): a geleira Pedersen em 1990 e em 2016, fotografias de M.~Tetreau e D.~Kurtz, U.S. National Park Service; a borda de uma plataforma de gelo ao largo da Península Antártica, NASA, Operation IceBridge (M.-J. Viñas); todas em domínio público (via Wikimedia Commons). O gráfico do gelo marinho usa o NSIDC Sea Ice Index, versão~4 (F.~Fetterer e colaboradores, 2025, doi:10.7265/a98x-0f50).
  \item Cap.~\ref{ch:g9:climate-past-and-present} (o clima da Terra, passado e presente): um trecho do testemunho de gelo GISP2 da Groenlândia, NOAA Paleoclimatology, domínio público (via Wikimedia Commons). Os registros vêm do NASA GISTEMP versão~4, do registro do gás carbônico de Mauna Loa (NOAA Global Monitoring Laboratory e Scripps Institution of Oceanography), da temperatura de EPICA Dome~C (J.~Jouzel e colaboradores, 2007), da composição do gás carbônico dos testemunhos de gelo da Antártida (B.~Bereiter e colaboradores, 2015) e do nível do mar do NOAA Laboratory for Satellite Altimetry; as afirmações da avaliação são as do Sexto Relatório de Avaliação do IPCC, Grupo de Trabalho~I, Sumário para Formuladores de Políticas (2021).
  \item Cap.~\ref{ch:g9:natural-risks} (conviver com os riscos naturais): o mapa de ameaças do Nevado del Ruiz, de T.~L. Wright e T.~C. Pierson, \emph{Living with Volcanoes}, U.S. Geological Survey Circular 1073 (1992); uma boia de alerta de tsunami inspecionada pela tripulação do USCGC \emph{William Hart}, U.S. Coast Guard, 2021; ambos em domínio público (via Wikimedia Commons).
  \item Cap.~\ref{ch:g9:earth-resources} (os recursos da Terra): a mina de Chuquicamata vista pelo ASTER, NASA/METI/AIST/Japan Space Systems e U.S./Japan ASTER Science Team, 2007, domínio público; uma camada de carvão num corte de estrada no condado de Belmont, em Ohio, fotografia de James St.~John, CC BY 2.0 (ambas via Wikimedia Commons). Os fatores de concentração são de D.~Perkins e colaboradores, \emph{Mineralogy}, e de S.~Earle, \emph{Environmental Geology}.
  \item Cap.~\ref{ch:g9:landforms} (as formas de relevo): as Seven Sisters, fotografia de Arild Vågen, 2017, CC BY-SA 4.0; dunas em crescente na costa do sul do Brasil, fotografia de um astronauta da NASA tirada da Estação Espacial Internacional, 2017, domínio público; estalactites na Carlsbad Cavern, fotografia do National Park Service por Daniel Leifheit, 2018, domínio público (todas via Wikimedia Commons).
  \item Cap.~\ref{ch:g9:maps-of-the-earth} (mapas da Terra inteira): o delta do rio Lena, imagem Landsat de 2000, USGS EROS Data Center via NASA Visible Earth, domínio público; um satélite GPS, concepção artística da NASA, domínio público (ambos via Wikimedia Commons). Litorais dos mapas-múndi do Natural Earth (domínio público).
\end{itemize}
